% chapters/05/contoh.tex
\section{Dua Contoh Terapan}
\label{sec:contoh-bab-5}

\subsection{Contoh Pertama: Menghitung Kecepatan Diagonal}

\begin{contohbox}
Seorang pemain menekan tombol kanan dan tombol bawah secara bersamaan.
Program menyusun vektor arah bernilai $(1, 1)$ tanpa menormalkannya, lalu
mengalikannya dengan kecepatan 200 piksel per detik.

Panjang vektor arah tersebut adalah

\[
|(1, 1)| = \sqrt{1^2 + 1^2} = \sqrt{2} \approx 1{,}414 .
\]

Karena kecepatan adalah arah dikalikan besarnya, maka kecepatan yang
sesungguhnya bekerja adalah

\[
200 \times 1{,}414 = 282{,}8 \text{ piksel per detik}.
\]

Objek bergerak 41 persen lebih cepat daripada ketika berjalan lurus, tanpa
sepengetahuannya dan tanpa pesan kesalahan apa pun dari program.
\end{contohbox}

Apabila vektor arah dinormalkan lebih dahulu, panjangnya menjadi satu, dan
kecepatannya kembali menjadi 200 piksel per detik untuk segala arah. Perbedaan
keduanya diperlihatkan pada Gambar~\ref{fig:normalisasi-diagonal}.

Kecepatan yang lebih besar itu membawa akibat yang lebih jauh
daripada sekadar terasa cepat. Jangkauan lompatan, jarak yang ditempuh sebelum
waktu habis, dan peluang menghindari musuh semuanya berubah ketika pemain
menempuh arah diagonal. Pada permainan yang mengandalkan ketepatan, keunggulan
41 persen itu cukup untuk membuat permainan tidak lagi seimbang. Karena itu,
penormalan vektor arah bukan kerapian tambahan, melainkan bagian dari
kebenaran perhitungan.

\subsection{Contoh Kedua: Mengapa Posisi Tidak Boleh Berupa Bilangan Bulat}

\begin{contohbox}
Sebuah objek bergerak dengan kecepatan 20 piksel per detik pada laju 60
bingkai per detik. Perpindahan setiap bingkainya adalah

\[
20 \div 60 = 0{,}333 \text{ piksel}.
\]

Misalkan posisi disimpan sebagai bilangan bulat. Pada setiap bingkai, program
menambahkan 0,333 lalu memotongnya menjadi bilangan bulat. Hasilnya
diperlihatkan pada Tabel~\ref{tab:posisi-bulat}.
\end{contohbox}

\begin{table}[htbp]
\centering
\caption{Posisi objek pada setiap bingkai menurut cara penyimpanannya}
\label{tab:posisi-bulat}
{\setlength{\tabcolsep}{4pt}\footnotesize%
\begin{tabularx}{\textwidth}{@{}C{1.6cm} C{3.4cm} Y@{}}
\toprule
\headrow \textbf{Bingkai} & \textbf{Posisi bila disimpan sebagai bilangan bulat} & \textbf{Posisi bila disimpan sebagai pecahan} \\
\midrule
1 & 0 & 0,333 \\
2 & 0 & 0,667 \\
3 & 0 & 1,000 \\
4 & 0 & 1,333 \\
5 & 0 & 1,667 \\
6 & 0 & 2,000 \\
\bottomrule
\end{tabularx}}
\tabsource{Kecepatan 20 piksel per detik pada 60 bingkai per detik. Pada kolom
kedua, nilai 0,333 terpotong menjadi 0 pada setiap bingkai sehingga nilainya
tidak pernah terakumulasi; pada kolom ketiga, nilai itu tersimpan dan
terakumulasi.}
\end{table}

Objek pada kolom kedua tidak akan bergerak sama sekali, berapa lama pun
program dijalankan. Nilai pecahan yang seharusnya terakumulasi terpotong
menjadi nol pada setiap bingkai, dan nilai yang terpotong itu hilang untuk
selamanya.

Gejala ini tidak muncul pada kecepatan tinggi. Pada kecepatan 200 piksel per
detik, perpindahan setiap bingkainya 3,33 piksel sehingga pemotongan hanya
memperlambat objek sedikit, dan kekeliruannya tidak terlihat. Kekeliruan
semacam ini karena itu berbahaya: program tampak bekerja pada pengujian
pertama, lalu gagal ketika Anda menurunkan kecepatan objek untuk keperluan
permainan yang sesungguhnya.

Kesimpulannya, simpanlah posisi sebagai pecahan dan gunakan \code{Rect} hanya
untuk menggambar serta memeriksa tabrakan. Pembagian tugas itu akan Anda pakai
terus sampai Bab~XVI.
